\section{Setup and notation}
\label{sec:setup}

\paragraph{Data.}
Comparisons $X=(x,y_A,y_B)$ are drawn from $P_X$. Before anything else, the data are split by
prompt into four parts (Section~\ref{sec:experiments}):
\begin{itemize}
  \item an untouched \emph{test set} $\mathcal T$, used for no acquisition decision, model choice or
        policy training;
  \item a \emph{development set} $\mathcal V$, used only for shared policy-step hyperparameters;
  \item a \emph{diagnostic set} $\mathcal G$, used only for E2;
  \item the \emph{pool} $X_1,\dots,X_N$.
\end{itemize}
A frozen judge is queried on every comparison, twice under protocol J1 (once per presentation
order). The dataset supplied for this project already contains a
human label $Y_i$ for every comparison. During training, only the pool labels indexed by
$\cI_n\subset[N]$, $|\cI_n|=n\ll N$, are revealed. \textbf{No new human annotation is collected.}

\paragraph{Human outcome and its embedding.}
Let $\cS_H$ be the human outcome space and $\QH^{*}(\cdot\mid X)$ the target conditional distribution.
We represent distributions through a feature map $\phi:\cS_H\to\cH$ with $\norm{\phi}\le1$, and write
$\cM$ for the closed convex hull of $\phi(\cS_H)$. The conditional mean embedding
$\E\{\phi(Y)\mid\cdot\}$ then always lies in $\cM$.

\begin{center}
\small
\renewcommand{\arraystretch}{1.15}
\begin{tabular}{@{}llll@{}}
\toprule
Human label & $\phi(y)$ & $\E\{\phi(Y)\mid\cdot\}$ & $\norm{\cdot}^2$ risk \\
\midrule
Binary $\{0,1\}$ & $y$ & $p_H^{*}$ & Brier / MSE \\
Categorical $\{1..K\}$ & one-hot $e_y$ & probability vector & multiclass Brier \\
Ordinal $\{1..K\}$ & $(K-1)^{-1/2}(\ind\{y\le k\})_{k<K}$ & scaled CDF vector & RPS (discrete CRPS) \\
General & kernel feature, $k(y,y)\le1$ & conditional mean embedding & conditional MMD$^2$ \\
\bottomrule
\end{tabular}
\end{center}

\noindent The ordinal scaling keeps $\norm{\phi}\le1$; the unscaled CDF vector has norm up to $\sqrt{K-1}$. For the kernel row see \citet{song2009hilbert,muandet2017kernel}.

Suppose $r$ annotators label a comparison independently given $X$. Their averaged embedding
$\bar\phi$ is then unbiased for $\mu_X:=\E\{\phi(Y)\mid X\}$. At the level of the localization
variable $Z$ introduced below,
$\mathrm{Var}(\bar\phi\mid Z)=r^{-1}\E\{\mathrm{Var}(\phi(Y)\mid X)\mid Z\}+\mathrm{Var}(\mu_X\mid Z)$.
Averaging removes only the within-comparison part of the noise.

\paragraph{Judge protocol.}
If the judge only returns a binary verdict, then $\QJ=\Ber(p_J)$ is determined by $p_J$. Its entropy,
variance and the empirical dispersion of repeated binary verdicts are all functions of $p_J$ (or of
its empirical estimate), so a ``distributional'' method has nothing extra to use. The distributional
question is meaningful only for a richer protocol. We therefore commit to the following one.
\begin{itemize}
  \item \textbf{J1 (primary): graded, logit-elicited, both orders.} The judge chooses among five
        labels $\{A\gg B, A>B, \text{tie}, B>A, B\gg A\}$. $\QJ^{AB}$ is the softmax over these
        label tokens with $A$ shown first; $\QJ^{BA}$ is the same with the order reversed, mapped
        back to $A/B$ labels. $\QJ=(\QJ^{AB},\QJ^{BA})$ has eight free coordinates, and its values
        are continuous.
  \item \textbf{J2 (secondary): sampled.} The same five-label prompt is sampled $M$ times in each
        order, giving empirical frequencies on a grid of step $1/M$. J2 is used only in the
        judge-budget experiment (E5). E5 is purely empirical: the summaries below do not determine
        the sampling law of their $M$-sample estimates (for example, the variance of $\hat t_J$
        depends on the two order-specific tie masses separately), so
        Proposition~\ref{prop:judge-budget} does not apply to them.
\end{itemize}

\paragraph{Judge prediction and link.}
A known link $T$ maps judge distributions into $\cM$ and defines the judge's prediction
$\mJ(X)=T\{\QJ(\cdot\mid X)\}$. For a binary human label we use the order-averaged win probability
$p_J=\tfrac12\{w(\QJ^{AB})+w(\QJ^{BA})\}$, where $w(Q)=Q(A\gg B)+Q(A>B)+\tfrac12Q(\text{tie})$. If the
human label is itself graded on the same scale, $T$ is the order-averaged embedding
$\tfrac12\{\E_{\QJ^{AB}}\phi+\E_{\QJ^{BA}}\phi\}$.

\paragraph{Localization variable.}
Kernel smoothing needs a low-dimensional covariate, and the raw text $X$ is not one. We localize in
\[
  Z=\zeta(X)=\bigl(\psi\{\QJ(\cdot\mid X)\},\,\xi(X)\bigr).
\]
Here $\psi$ collects judge summaries that are \emph{not} functions of $p_J$:
\begin{itemize}
  \item tie mass $t_J=\tfrac12\{\QJ^{AB}(\text{tie})+\QJ^{BA}(\text{tie})\}$;
  \item strength $s_J$, the order-averaged mass on $\{A\gg B,B\gg A\}$;
  \item order inconsistency $o_J=|w(\QJ^{AB})-w(\QJ^{BA})|$.
\end{itemize}
$\xi$ holds optional content features. Discrete features such as task category are handled by
stratification: \dfsp{} is fitted separately within each stratum, so the kernel only acts on
continuous coordinates. The scalar and distributional variants always use the same $\xi$:
\[
  Z_s=(p_J,\xi),\qquad Z_d=(p_J,t_J,o_J,\xi)\ \text{(default; $s_J$ is an ablation)}.
\]
\paragraph{Scope: binary human labels.}
The main method targets binary human labels. Its offset is $\mJ=p_J$, which is a coordinate of
$Z$, so $\mJ(z)$ is well defined. For graded human labels this fails: the graded offset depends on
the full order-averaged five-label vector, and $(p_J,t_J,o_J)$ does not determine that vector. (Two
vectors can share $p_J=.5$, $t_J=.2$, $o_J=0$ and even the same strength mass, yet have different
CDFs.) Graded human targets are therefore an extension, which would need $Z$ to contain the whole
vector ($d=4$ after the simplex constraint). They are not part of the main plan.

\paragraph{Feasible support.}
The coordinates are deterministically constrained. Write $w_1=w(\QJ^{AB})$ and $w_2=w(\QJ^{BA})$,
so $p_J=(w_1+w_2)/2$ and $o_J=|w_1-w_2|$. Then:
\begin{itemize}
  \item $w_i\in[0,1]$ gives $o_J\le2\min(p_J,1-p_J)$;
  \item $w(Q)\in[\tfrac12Q(\text{tie}),1-\tfrac12Q(\text{tie})]$ gives $t_J\le2\min(p_J,1-p_J)$ and
        $t_J+o_J\le1$.
\end{itemize}
For unrestricted pairs of five-label distributions these are also sufficient. (Take
$w_{1,2}=p\pm o/2$ and split the tie mass subject to $t_i\le2\min(w_i,1-w_i)$.) The feasible region
is therefore exactly
\[
  \cZ=\bigl\{(p,t,o):\ 0\le t\le2\min(p,1-p),\ 0\le o\le\min\{2\min(p,1-p),\,1-t\}\bigr\}
  \subsetneq[0,1]^3 ,
\]
a polytope with nonempty interior. The theory of Section~\ref{sec:theory} is stated on a regular
subset $\cZ_0\subseteq\cZ$, specified after the pilot, on which the densities are assumed bounded
below. It is not stated on the cube. Density estimation, occupancy and acquisition use the same
$\cZ_0$. Under J1 the coordinates are continuous. The E0 feasibility pilot is a
\emph{diagnostic} for support and degeneracy (for example, $t_J$ concentrating near $0$); it cannot
verify a population density lower bound.

\paragraph{Estimand.}
\[
  \muH(z)=\E\{\phi(Y)\mid Z=z\}\in\cM .
\]
For binary labels, $\muH=p_H^{*}$. The estimand depends on $\zeta$: it is $\mu_X$ coarsened to $Z$,
and the coarsening gap $B_\zeta=\E\norm{\muH(Z)-\mu_X}^2$ is irreducible for any method that uses
only $Z$ (Section~\ref{sec:scalar-vs-dist}).

\keybox{\textbf{\textcolor{BrickRed}{Design decision D1: the localization variable $\zeta$.}}
Default: $Z_d=(p_J,t_J,o_J)$ on a regular subset $\cZ_0$ of the polytope $\cZ$, stratified by task
category, so $d=3$. A
high-dimensional content embedding is excluded from the main method because of the curse of
dimensionality. The representation is frozen after the E0 feasibility pilot, using judge outputs
only and no human labels. If $t_J$ or $o_J$ is nearly degenerate there, it is dropped at that
point.}
